% The four panels share the content and two-by-two layout of the source figure.
\begin{subfigure}[t]{0.49\textwidth}
  \centering
  \resizebox{\linewidth}{!}{%
    \begin{tikzpicture}[>=Stealth, every node/.style={font=\scriptsize}]
      \useasboundingbox (0,-0.75) rectangle (6,2.2);
      \node (top-left) at (2.05,1.82)
        {$\tau^{*}\!\left(AB^{n}\dfrac{R}{R^{\times}}\right)$};
      \node (top-right) at (4.75,1.82)
        {$B^{n}\dfrac{R}{R^{\times}}$};
      \node (left) at (0.55,0.35) {$X$};
      \node (middle) at (2.8,0.35) {$X$};
      \node (right) at (5.0,0.35) {$\boldsymbol{B}R^{\times}$};
      \draw[->] (top-left) -- (top-right);
      \draw[->] (top-left.south) -- (middle.north);
      \draw[->] (top-right.south) -- (right.north);
      \draw[dashed,->] (left.north east) -- (top-left.south west);
      \draw[double,double distance=1.2pt] (left.east) -- (middle.west);
      \draw[->] (middle) -- node[above] {$\tau$} (right);
      \draw[->,blue,thick] (right.south) -- (5.0,-0.55) --
        (0.55,-0.55) -- (left.south);
      \node[blue,fill=white,inner sep=1pt] at (2.78,-0.55) {$\Gamma_R^{*}$};
    \end{tikzpicture}%
  }
  \caption{交换关系}
\end{subfigure}\hfill
\begin{subfigure}[t]{0.49\textwidth}
  \centering
  \resizebox{\linewidth}{!}{%
    \begin{tikzpicture}[>=Stealth, every node/.style={font=\scriptsize}]
      \useasboundingbox (0,-0.75) rectangle (6,2.2);
      \node (id) at (0.45,1.72) {$\mathrm{id}_S\in$};
      \node (hom-top) at (1.95,1.72) {$\operatorname{Hom}_{\mathcal C}(S,S)$};
      \node (a-top) at (4.00,1.72) {$A(S)$};
      \node (u) at (5.55,1.72) {$\ni u_S$};
      \node (f) at (0.45,0.25) {$f\in$};
      \node (hom-bottom) at (1.95,0.25) {$\operatorname{Hom}_{\mathcal C}(T,S)$};
      \node (a-bottom) at (4.00,0.25) {$A(T)$};
      \node (phi) at (5.55,0.25) {$\ni\phi_T(f)$};
      \draw[->] (hom-top) -- node[above] {$\phi_S$} (a-top);
      \draw[->] (hom-bottom) -- node[above] {$\phi_T$} (a-bottom);
      \draw[->] (id) -- (f);
      \draw[->] (hom-top) -- node[right] {$f^{*}$} (hom-bottom);
      \draw[->] (a-top) -- node[right] {$A(f)$} (a-bottom);
      \draw[->] (u) -- (phi);
    \end{tikzpicture}%
  }
  \caption{范畴关系}
\end{subfigure}

\par\smallskip

\begin{subfigure}[t]{0.49\textwidth}
  \centering
  \resizebox{\linewidth}{!}{%
    \begin{tikzpicture}[>=Stealth, every node/.style={font=\scriptsize},
      state/.style={draw,circle,minimum size=12mm,inner sep=1pt,fill=blue!12}]
      \useasboundingbox (0,-0.75) rectangle (6,2.2);
      \node[state] (reading) at (1.45,0.45) {reading};
      \node[state] (eof) at (3.35,0.45) {eof};
      \node[state,double,double distance=2pt] (closed) at (5.2,0.45) {closed};
      \draw[->] (0.05,0.45) -- node[above] {open(path)} (reading.west);
      \draw[->] (reading.east) -- node[above] {read()} (eof.west);
      \draw[->] (eof.east) -- node[above] {close()} (closed.west);
      \draw[->] (reading) to[loop above] node[above=2pt] {read()} (reading);
      \draw[->] (reading.south east) to[bend right=32]
        node[below] {close()} (closed.south west);
    \end{tikzpicture}%
  }
  \caption{离散状态机}
\end{subfigure}\hfill
\begin{subfigure}[t]{0.49\textwidth}
  \centering
  \resizebox{\linewidth}{!}{%
    \begin{tikzpicture}[>=Stealth, every node/.style={font=\scriptsize},
      tree-node/.style={draw=gray!55,fill=white,align=center,font=\tiny,
        text width=13mm,minimum height=6mm,inner sep=1.5pt},
      tree-node-wide/.style={tree-node,text width=18mm}]
      \useasboundingbox (0,-0.75) rectangle (6,2.2);
      \node[tree-node] (root) at (3.0,1.95) {Entrate\\pubblicitarie};
      \node[tree-node] (impr) at (1.55,0.85) {Numero di\\impression};
      \node[tree-node-wide] (gain) at (4.45,0.85) {Guadagno medio\\per impression};
      \node[tree-node] (traffic) at (0.55,-0.28) {Traffico\\sul sito};
      \node[tree-node] (ratio) at (2.25,-0.28) {Impression/\\visitatori};
      \node[tree-node] (sales) at (3.75,-0.28) {Modalità\\di vendita};
      \node[tree-node] (placement) at (5.45,-0.28) {Tipo di\\posi\-zionamento};
      \draw[->] (root.south west) -- (impr.north);
      \draw[->] (root.south east) -- (gain.north);
      \draw[->] (impr.south west) -- (traffic.north);
      \draw[->] (impr.south east) -- (ratio.north);
      \draw[->] (gain.south west) -- (sales.north);
      \draw[->] (gain.south east) -- (placement.north);
    \end{tikzpicture}%
  }
  \caption{层级状态树}
\end{subfigure}
